\documentclass{article}
\input{../../../lib.sty}
\title{\texttt{fn sample\_negative\_binomial\_integer}}
\author{Michael Shoemate}
\begin{document}
\maketitle

This document proves that the implementation of \texttt{sample\_negative\_binomial\_integer}
in \texttt{mod.rs} satisfies its proof definition.

\section{Hoare Triple}
\subsection*{Precondition}
$\texttt{shape} \in \mathbb{N}_{>0}$ and $\texttt{x} \in \mathbb{Q}_{>0}$.

\subsection*{Pseudocode}
\lstinputlisting[language=Python,firstline=2,escapechar=|]{./pseudocode/sample_negative_binomial_integer.py}

\subsection*{Postcondition}
\label{postcondition}
For any setting of the input parameters \texttt{shape} and \texttt{x} such that the given preconditions hold,
\texttt{sample\_negative\_binomial\_integer} either returns \texttt{Err(e)} due to a lack of system entropy,
or \texttt{Ok(out)}, where \texttt{out} is distributed as a negative binomial random variable on $\mathbb{N}_0$
with shape \texttt{shape} and success probability $1 - e^{-x}$, parameterized as the number of failures before
the \texttt{shape}-th success.

\section{Proof}
Assume the preconditions are met.

\begin{lemma}
\label{lem:err-e}
\texttt{sample\_negative\_binomial\_integer} only returns \texttt{Err(e)} when there is a lack of system entropy.
\end{lemma}
\begin{proof}
Under the stated preconditions, the preconditions of
\rustdoc{traits/samplers/cks20/fn}{sample\_geometric\_exp\_fast}
on line \ref{line:geom_draw} are met in every iteration of the loop.
By its proof definition, each invocation of
\rustdoc{traits/samplers/cks20/fn}{sample\_geometric\_exp\_fast}
either returns a geometrically distributed sample or \texttt{Err(e)} due to a lack of system entropy.
The function has no other fallible operations. Therefore any returned error must be due to a lack of system entropy.
\end{proof}

\begin{lemma}
\label{lem:geom-sum}
If the outcome of \texttt{sample\_negative\_binomial\_integer} is \texttt{Ok(out)}, then \texttt{out}
is distributed as a negative binomial random variable on $\mathbb{N}_0$ with shape \texttt{shape}
and success probability $1 - e^{-x}$, parameterized as the number of failures before the
\texttt{shape}-th success.
\end{lemma}
\begin{proof}
Let $m$ denote the value of \texttt{shape}.
For each of the $m$ iterations of the loop, line \ref{line:geom_draw} draws an independent sample
$G_i \sim Geometric(1 - e^{-x})$ on $\mathbb{N}_0$.
The algorithm returns
$K = \sum_{i=1}^m G_i$.

For any $k \in \mathbb{N}_0$, the event $\{K = k\}$ is equivalent to choosing a vector
$(g_1, \ldots, g_m) \in \mathbb{N}_0^m$ such that $\sum_{i=1}^m g_i = k$.
There are $\binom{k + m - 1}{k}$ such vectors.
Each vector has probability
\begin{equation*}
\prod_{i=1}^m (1 - e^{-x}) e^{-x g_i}
= (1 - e^{-x})^m e^{-xk}.
\end{equation*}
Therefore
\begin{equation*}
\Pr[K = k]
= \binom{k + m - 1}{k} (1 - e^{-x})^m e^{-xk}.
\end{equation*}
This is exactly the negative binomial distribution with shape $m$ and success probability $1 - e^{-x}$,
parameterized as the number of failures before the $m$-th success.
\end{proof}

\begin{proof}
\ref{postcondition} holds by Lemma~\ref{lem:err-e} and Lemma~\ref{lem:geom-sum}.
\end{proof}

\end{document}
